Una biblioteca va fer 300 préstecs en una setmana, entre llibres, còmics i pel·lícules. Els còmics i el
doble de les pel·lícules en van ser 80 més que els llibres. Si se suma el 20\,\% dels llibres, el 30\,\%
dels còmics i el 50\,\% de les pel·lícules, que són els que es van tornar amb retard, s'obtenen 88
préstecs.

\begin{apartats}

\apartat[1,25]{0,75}
Planteja un sistema d'equacions que permeti saber quants préstecs de cada tipus es van fer.

\begin{solucio}
Siguin $x$, $y$ i $z$ els préstecs de llibres, de còmics i de pel·lícules:
\[
\left.\begin{aligned}x+y+z&=300\\y+2z&=x+80\\0{,}2x+0{,}3y+0{,}5z&=88\end{aligned}\right\}
\ \Longleftrightarrow\
\left.\begin{aligned}x+y+z&=300\\-x+y+2z&=80\\2x+3y+5z&=880\end{aligned}\right\}
\]
(la tercera, multiplicada per 10).
\end{solucio}

\begin{tria}{resolucio}
\itemtria{resol}{1}{1,25}
Resol el sistema pel mètode de Gauss i interpreta'n la solució.

\begin{solucio}
\[
\left(\begin{array}{ccc|c}1&1&1&300\\-1&1&2&80\\2&3&5&880\end{array}\right)
\xrightarrow[F_3-2F_1]{F_2+F_1}
\left(\begin{array}{ccc|c}1&1&1&300\\0&2&3&380\\0&1&3&280\end{array}\right)
\xrightarrow{2F_3-F_2}
\left(\begin{array}{ccc|c}1&1&1&300\\0&2&3&380\\0&0&3&180\end{array}\right).
\]
$z=60$; $2y=380-180$, $y=100$; $x=300-100-60=140$. Es van fer 140 préstecs de llibres, 100 de còmics
i 60 de pel·lícules.
\end{solucio}

\itemtria{model-economic}{1}{1,25}
Aquest altre problema és d'un model econòmic. Una empresa elèctrica i una d'aigua es proveeixen
mútuament: per produir 1\,€ d'electricitat es consumeixen 0,10\,€ d'electricitat i 0,20\,€ d'aigua, i
per produir 1\,€ d'aigua, 0,30\,€ d'electricitat i 0,10\,€ d'aigua. A més, fora de les dues empreses es
demanen 60\,€ d'electricitat i 30\,€ d'aigua. Quant ha de produir cada empresa perquè la producció
cobreixi el consum de les dues empreses i la demanda de fora?

\begin{solucio}
Siguin $E$ i $W$ les produccions, en euros, d'electricitat i d'aigua. La producció de cada una ha de ser
igual al que en consumeixen les dues empreses més la demanda exterior:
\[
\left.\begin{aligned}E&=0{,}1E+0{,}3W+60\\W&=0{,}2E+0{,}1W+30\end{aligned}\right\}
\ \Longleftrightarrow\
\left.\begin{aligned}9E-3W&=600\\-2E+9W&=300\end{aligned}\right\}
\]
De la primera, $W=3E-200$; substituint, $-2E+27E-1800=300$, $E=84$, i $W=52$. L'empresa elèctrica ha
de produir 84\,€ i la d'aigua, 52\,€.
\end{solucio}
\end{tria}

\begin{nomesllarg}
\apartat{0,75}
Calcula el determinant de la matriu de coeficients del sistema de l'apartat a). Si les dades de la
setmana següent fossin unes altres, amb les mateixes relacions, el sistema tindria sempre una única
solució? I el problema? Raona-ho.

\begin{solucio}
$\begin{vmatrix}1&1&1\\-1&1&2\\2&3&5\end{vmatrix}=5+4-3-2-6+5=3\ne0$. Com que el determinant no depèn dels
termes independents, el sistema és compatible determinat amb qualsevol dada: té sempre una única
solució.\\
El problema, però, només té solució si aquesta solució és de nombres naturals: amb unes altres dades,
podria sortir un nombre negatiu o no enter, que no correspondria a cap nombre de préstecs.
\end{solucio}
\end{nomesllarg}

\end{apartats}
